\documentclass[11pt,a4paper]{article}
% Shared preamble for RuPhO English problem statements (match T1 2026 style).
% Use: \input{latex/rupho_en_preamble.tex} after \documentclass.
\usepackage[margin=1in]{geometry}
\usepackage{amsmath,amssymb}
\usepackage{graphicx}
\usepackage{float}
\usepackage{enumitem}
\usepackage{microtype}
\usepackage{caption}
\usepackage{tabularx}
\usepackage{hyperref}

\captionsetup{font=small,labelfont=bf,skip=6pt}

\setlist[itemize]{leftmargin=1.2cm}
\setlist[enumerate]{leftmargin=1.2cm}

% One outer box per subproblem: [id | statement | points] (no inner rules)
\newcommand{\problembox}[3]{%
  \par\addvspace{0.42em}%
  \noindent
  \setlength{\fboxsep}{7.5pt}%
  \setlength{\fboxrule}{0.95pt}%
  \fbox{%
    \begin{minipage}{\dimexpr\linewidth-2\fboxsep-2\fboxrule\relax}
    \begin{tabularx}{\linewidth}{@{}>{\raggedright\arraybackslash\bfseries}p{2.1em} @{\hspace{0.9em}} >{\raggedright\arraybackslash}X @{\hspace{0.9em}} >{\raggedleft\arraybackslash\bfseries}p{2.4em}@{}}
    #1 & #2 & #3
    \end{tabularx}
    \end{minipage}%
  }%
  \par\addvspace{0.32em}%
}

\title{\textbf{Composite capacitor}}
\author{\normalsize X16 (2016) --- English translation}
\date{}
\begin{document}
\maketitle

\section*{Capacitor charge}

Between plates 1 and 3 of a flat capacitor, a thin metal plate 2 with a charge $Q$ was placed parallel to the plates of the capacitor (Fig. 1). A capacitor is connected to a battery with an emf equal to $U_0$. The areas of all three plates are the same and equal to $S$, and the distance between them is much less than $\sqrt{S}$.

\begin{figure}[H]
    \centering
    \includegraphics[width=0.72\textwidth]{T6_EN_assets/fig1.jpg}
    \caption{Figure 1.}
\end{figure}

\problembox{A1}{Find the charges on plates 1 and 3 of the capacitor depending on the distance $x$ between plates 1 and 2.}{1.00}

\problembox{A2}{How much work $A$ must be done to move plate 2 from plate 1 to plate 3?}{1.00}

Plate 2 was discharged, and the volumes between the plates were filled with dielectric liquids (Fig. 2) with the same dielectric constant $\varepsilon$ but with different resistivities $\rho_1$ and $\rho_2$ ($\rho_1 < \rho_2$).

\begin{figure}[H]
    \centering
    \includegraphics[width=0.72\textwidth]{T6_EN_assets/fig2.jpg}
    \caption{Figure 2.}
\end{figure}

\problembox{A3}{Find the magnitude and direction of the force acting on plate 2 from the electric field.}{3.00}

\vspace{1em}
\noindent\textit{Source:} \url{https://pho.rs/p/4004}

\end{document}